% ==============================================================
%  Chapter 2: AI for Developers
% ==============================================================
\chapter{AI for Developers}
\label{ch:developers}

AI has changed software development more than any tool since the internet.
Used well, it can double productivity on repetitive tasks, compress learning curves, and
eliminate the blank-page problem.
Used carelessly, it produces plausible-looking code that doesn't work, introduces subtle
bugs, and erodes the developer's own skills over time.
This chapter is about using it well.

\section{Code Generation}
\label{sec:codegen}

AI excels at generating code for well-understood patterns: CRUD operations, boilerplate
configuration, script scaffolding, API calls to documented libraries, and data
transformations.
These are tasks the model has seen thousands of times across millions of repositories.

\textbf{How to generate code that actually works:}
\begin{enumerate}
  \item State language, version, and framework first.
  \item Describe the input and desired output precisely.
  \item List constraints (no external libraries, max lines, error handling style).
  \item Ask for an explanation alongside the code.
  \item Run it before trusting it.
\end{enumerate}

\begin{tipbox}[Good Generation Prompt Structure]
\textit{``Write a Python 3.13 function that accepts a list of integers and returns a
dict mapping each integer to its square root, rounded to 2 decimal places. Skip
negative values silently. No external libraries. Include a docstring and one usage
example.''}
\end{tipbox}

\textbf{Where generation fails:}
\begin{itemize}
  \item Business-specific logic that does not appear in public training data
  \item Code that requires knowledge of your internal data models or APIs
  \item Novel algorithm design with non-obvious tradeoffs
  \item Accurate use of APIs that changed after the model's training cutoff
\end{itemize}

\begin{warnbox}
Always verify library function signatures against official documentation.
AI frequently generates plausible-looking but non-existent method names, especially for
less popular libraries or recent versions.
\end{warnbox}

\section{Debugging with AI}
\label{sec:debugging}

The most common mistake in AI-assisted debugging is pasting only the error message.
AI cannot know what you expected to happen, what the wider system state is, or what
you have already tried.

\begin{infobox}[The Complete Debugging Prompt]
Give AI all five elements:
\begin{enumerate}
  \item \textbf{Full error message} --- complete stack trace, not the last line only
  \item \textbf{Relevant code} --- the failing function plus any callers or dependencies
  \item \textbf{Expected behaviour} --- what should happen
  \item \textbf{Actual behaviour} --- what is happening
  \item \textbf{What you've already tried}
\end{enumerate}
Ask: \textit{``What are the three most likely root causes and how would I confirm each?''}
\end{infobox}

This approach prompts AI to reason diagnostically rather than guess a fix.
You then test each hypothesis rather than blindly applying the first suggestion.

AI is particularly strong at:
\begin{itemize}
  \item Decoding cryptic compiler and runtime error messages
  \item Identifying off-by-one errors and null-pointer patterns
  \item Explaining what a stack trace is actually telling you
  \item Suggesting diagnostic logging or breakpoint placement strategies
\end{itemize}

\section{Code Review and Refactoring}
\label{sec:review}

AI makes a useful first-pass reviewer.
It catches common issues quickly: missing null checks, inefficient loops, inconsistent
error handling, magic numbers, and style violations.
It is \emph{not} a replacement for human review of business logic or security-critical
paths.

\textbf{Effective review prompts:}
\begin{itemize}
  \item \textit{``Review this code for correctness, readability, and edge cases. List issues by severity.''}
  \item \textit{``What design pattern would improve this structure and why?''}
  \item \textit{``Refactor this to be more readable. Explain every change you make.''}
  \item \textit{``What would break if the input to this function is null, empty, or very large?''}
  \item \textit{``Does this code have any performance bottlenecks? What is the time complexity?''}
\end{itemize}

\begin{warnbox}
Never skip human code review because AI said the code looks fine.
AI regularly misses business-logic errors, context-dependent bugs, and subtle security
flaws. Use AI as a first pass, not a final gate.
\end{warnbox}

\section{Testing}
\label{sec:testing}

AI generates test structure and happy-path cases quickly.
It tends to miss edge cases, negative paths, and integration scenarios unless you
specifically ask for them.

\textbf{A reliable test-generation workflow:}
\begin{enumerate}
  \item Paste the function under test with its docstring.
  \item Ask AI to \emph{list} the cases to test before generating any code.
    (\textit{``What are all the cases I should test for this function?''})
  \item Review the list --- add anything missing.
  \item Ask AI to generate tests for the full list using your testing framework.
  \item Review the tests: do they actually exercise the right behaviour?
    Would they catch a real regression?
\end{enumerate}

\begin{tipbox}
Ask for the test list before the test code.
This catches coverage gaps before implementation and produces measurably better tests
than asking for code directly.
\end{tipbox}

\section{Documentation}
\label{sec:documentation}

Documentation is one of AI's strongest practical contributions.
It reads existing code accurately and produces docstrings, README sections, API
references, and Architecture Decision Records quickly.

Useful prompts:
\begin{itemize}
  \item \textit{``Write a docstring for this function: parameters, return value, exceptions thrown, and one usage example.''}
  \item \textit{``Generate a README for this module: what it does, install steps, configuration, and two usage examples.''}
  \item \textit{``Write an Architecture Decision Record for the choice to use X over Y, given these constraints: [list].''}
  \item \textit{``Summarise this git diff as a commit message in Conventional Commits format.''}
\end{itemize}

\section{Learning New Technologies}
\label{sec:learning}

AI is the fastest onboarding tool available for new languages, frameworks, or tools.
It explains concepts, provides runnable examples, and answers follow-up questions
immediately --- without you waiting for a colleague or reading 200-page documentation.

\textbf{Effective learning patterns:}
\begin{itemize}
  \item Bridge from what you know:
    \textit{``I'm an experienced Python developer. Explain Go goroutines using Python
    async/await as a reference point.''}
  \item Ask for the pitfalls first:
    \textit{``What are the top 5 mistakes a Python developer makes when first learning Go?''}
  \item Learn by doing:
    \textit{``Give me a working example of channels in Go, then give me an exercise to
    test my understanding.''}
\end{itemize}

\begin{warnbox}[Skill Maintenance]
The risk of AI-assisted learning is passive consumption.
Deliberately write code without AI assistance regularly.
If you cannot explain or modify AI-generated code you have committed to your codebase,
you have a skills debt that will compound.
\end{warnbox}

\section{Security in AI-Generated Code}
\label{sec:security}

Security is the highest-risk area of AI-assisted development.
AI knows security patterns but applies them inconsistently.
It can generate code with SQL injection, broken authentication, or missing input
validation without any warning, and will not flag the issue unless you ask.

\begin{warnbox}[Security Non-Negotiables]
\begin{itemize}
  \item Manually review all authentication, authorisation, and data-handling code
  \item Run a static security analyser on AI-generated code (Semgrep, Bandit, ESLint
        security plugins)
  \item Never use AI-generated cryptographic implementations without expert review
  \item Treat AI-generated input validation as a starting point only
  \item Never commit code containing secrets or credentials, regardless of origin
\end{itemize}
\end{warnbox}

\section{Practical Prompt Templates for Developers}
\label{sec:dev-templates}

\begin{table}[htbp]
\centering
\caption{Proven prompt templates by development task}
\begin{tabularx}{\linewidth}{l X}
\toprule
\textbf{Task} & \textbf{Prompt Pattern} \\
\midrule
Understand code     & ``Explain this code line by line, then summarise what business
                       problem it solves.'' \\[4pt]
Generate boilerplate & ``Write a [type] in [lang/framework] that does [X].
                       Constraints: [Y]. Include a docstring.'' \\[4pt]
Debug               & ``Given this code and error, what are 3 possible root causes?
                       How do I confirm each?'' \\[4pt]
Refactor            & ``Refactor this for readability using [principle/pattern].
                       Explain each change.'' \\[4pt]
Write tests         & ``List test cases for this function, then generate them
                       in [framework]. Include edge cases.'' \\[4pt]
Learn a concept     & ``I know [X]. Explain [Y] using [X] as a reference.
                       Give a concrete runnable example.'' \\[4pt]
Translate code      & ``Convert this [lang A] to [lang B]. Note any
                       semantic differences.'' \\[4pt]
Security review     & ``Review this for OWASP Top 10 issues.
                       List each issue and how to fix it.'' \\
\bottomrule
\end{tabularx}
\label{tab:dev-prompts}
\end{table}
